\documentclass[12pt]{article}
\usepackage{pst-all}
\usepackage[french]{babel}
\usepackage{tikz}
\usepackage{tkz-tab}
\usepackage[T1]{fontenc}
\usepackage{eurosym}
\usepackage{amssymb}
\usepackage{float}
\usepackage{blkarray}
\usepackage{dsfont}
\usepackage{amsfonts}
\usepackage{amsmath}
\usepackage{eurosym}
\usepackage{array}
\usepackage{graphicx}
\parindent=0cm
\usepackage{geometry}
\usepackage{color}
\usepackage{graphicx}
\geometry{hmargin=1.5cm,vmargin=1.5cm}
\usepackage{amsmath, amssymb, amsfonts, amscd, amsthm}
\usepackage{pifont}
\usepackage{fancyhdr} 
\pagestyle{fancy}
\usepackage{multirow}
\renewcommand{\headrulewidth}{1pt}
\renewcommand{\footrulewidth}{1pt}
%\fancyhead[L]{Mr NOM}
\fancyhead[R]{ }
\fancyhead[C]{$5^e$}





\usepackage[utf8]{inputenc}
\usepackage{tcolorbox}% charge aussi pgf, verbatim, etoolbox et environ
\tcbuselibrary{theorems}% charge aussi amsmath
\usepackage{xcolor}

\definecolor{propcolor}{rgb}{.804,.064,.464}
\definecolor{defcolor}{rgb}{0,.545,0}
\definecolor{thcolor}{rgb}{0.932,0,0}
\definecolor{remcolor}{rgb}{1,0.5,0}
\definecolor{excolor}{rgb}{.094,.455,.804}

\usepackage{cleveref}

\newtcbtheorem[crefname={definition}{definitions}]{Definition}{Definition}{defstyle}{def}


\usepackage{varwidth} 
\tcbuselibrary{skins}

\newtcbtheorem[use counter from=Definition]{defin}{Définition}
{enhanced,frame empty, theorem name,interior empty,colframe=defcolor,coltitle=defcolor,fonttitle=\bfseries,
	colbacktitle=white,
	borderline={0.5mm}{0mm}{defcolor},
	%borderline={0.5mm}{0mm}{ForestGreen!50!white,dashed},
	attach boxed title to top left={yshift=-2mm, xshift = 2mm},
	boxed title style={boxrule=0.4pt}, varwidth boxed title}{theo}

\newtcbtheorem[use counter from=Definition]{theo}{Théorème}
{enhanced,frame empty,theorem name,interior empty,colframe=thcolor,coltitle=thcolor,fonttitle=\bfseries,
	colbacktitle=white,
	borderline={0.5mm}{0mm}{thcolor},
	%borderline={0.5mm}{0mm}{ForestGreen!50!white,dashed},
	attach boxed title to top left={yshift=-2mm, xshift = 2mm},
	boxed title style={boxrule=0.4pt}, varwidth boxed title}{theo}

\newtcbtheorem[use counter from=Definition]{prop}{Propriété}
{enhanced,frame empty,theorem name,interior empty,colframe=propcolor,coltitle=propcolor,fonttitle=\bfseries,
	colbacktitle=white,
	borderline={0.5mm}{0mm}{propcolor},
	%borderline={0.5mm}{0mm}{ForestGreen!50!white,dashed},
	attach boxed title to top left={yshift=-2mm, xshift = 2mm},
	boxed title style={boxrule=0.4pt}, varwidth boxed title}{theo}


\newtcbtheorem[use counter from=Definition]{ex}{Exemple}
{enhanced,frame empty,theorem name,interior empty,colframe=excolor,coltitle=excolor,fonttitle=\bfseries,
	colbacktitle=white,
	borderline={0.5mm}{0mm}{excolor},
	%borderline={0.5mm}{0mm}{ForestGreen!50!white,dashed},
	attach boxed title to top left={yshift=-2mm, xshift = 2mm},
	boxed title style={boxrule=0.4pt}, varwidth boxed title}{theo}

\newtcbtheorem[use counter from=Definition]{rem}{Remarque}
{enhanced,frame empty,theorem name,interior empty,colframe=remcolor,coltitle=remcolor,fonttitle=\bfseries,
	colbacktitle=white,
	borderline={0.5mm}{0mm}{remcolor},
	%borderline={0.5mm}{0mm}{ForestGreen!50!white,dashed},
	attach boxed title to top left={yshift=-2mm, xshift = 2mm},
	boxed title style={boxrule=0.4pt}, varwidth boxed title}{theo}


\newif\iftrou\troufalse
%mettre true à la place de false pour mettre des pointillés

\newlength{\trouwidth}

\newcommand{\trou}[1]{%
	\iftrou%
	\settowidth{\trouwidth}{#1}%
	\hbox to \trouwidth{\dotfill}%
	\else%
	#1%
	\fi}


\newcommand\pointille[1][1cm]{\makebox[#1]{\dotfill}}


\begin{document}


%%%%%%%%%%%%%%%%%%%%%%%%%%%%%%%%%%%%%%%%%%%%%%%%%%%%%%%%%%%%


\begin{center}
		\rule{\linewidth}{0.5mm} \\[0.4cm]
		{ \huge \bfseries Traitement de données\\[0.4cm] }
		\rule{\linewidth}{0.5mm} \\
\end{center}



%%%%%%%%%%%%%%%%%%%%%%%%%%%%%%%%%%%%%%%%%%%%%%%%%%%%%%%%%%%%%%%%%%%%


\section{Calculer des effectifs, des fréquences, une moyenne}

\begin{defin}{}{}
	Dans une série de données :\\
	- l'\textbf{effectif} d'une donnée est le nombre de fois où cette donnée apparaît.\\
	- l'\textbf{effectif total} est la somme de tous les effectifs.\\
	- la \textbf{fréquence} d'une donnée est le quotient de son effectif par l'effectif total:
	  \[ \text{fréquence d'une donnée} = \displaystyle\frac{\text{effectif de la donnée}}{\text{effectif total}}\]
\end{defin}


\begin{ex}{}{}
	Voici les réponses d'un groupe d'élèves à la question : "Quelle est votre couleur préférée ?" : \\
	\textcolor{blue}{bleu} - \textcolor{red}{rouge} - \textcolor{blue}{bleu} - \textcolor{olive}{vert} - \textcolor{purple}{violet} - \textcolor{blue}{bleu} - \textcolor{olive}{vert} -  \textcolor{red}{rouge} - \textcolor{olive}{vert} - \textcolor{olive}{vert} - \textcolor{purple}{violet} -\textcolor{purple}{violet} - \textcolor{pink}{rose} - \textcolor{olive}{vert} - \textcolor{orange}{orange} - \textcolor{blue}{bleu} - \textcolor{red}{rouge} - \textcolor{blue}{bleu} - \textcolor{orange}{orange} - \textcolor{olive}{vert}\\
	On peut regrouper cette série de données dans un tableau, appelé \textbf{tableau d'effectifs} :
	\begin{center}
				\includegraphics[width=0.85\linewidth]{don1}
	\end{center}
\end{ex}

\begin{rem}{}{}
	Les fréquences et les effectifs sont \trou{proportionnels}.\\
	Une fréquence peut être donnée sous forme de fraction, de nombre décimal ou de pourcentage; elle est comprise entre \trou{$0$} et \trou{$1$}.\\
	La somme de toutes les fréquences est égale à \trou{$1$}.
\end{rem}


\begin{defin}{}{}
	La \textbf{moyenne} d'une série de données est égale au quotient de la somme de ces données par l'effectif total:
	\[ \text{moyenne} = \displaystyle\frac{\text{somme des données}}{\text{effectif total}}\]
\end{defin}


\begin{rem}{}{}
	La moyenne n'est pas nécessairement égale à l'une des données.\\
	La moyenne est toujours comprise entre la plus petite et la plus grande valeur de la série. 
\end{rem}

\begin{ex}{}{}
	Voici les notes sur $20$ obtenues par Alice en mathématiques au premier trimestre: $11$ ; $12,5$ ; $14$ ; $9,5$ ; $13$ \\
	Pour calculer sa moyenne, on calcule la somme de ses notes, que l'on divise par le nombre de notes: $\displaystyle\frac{\trou{11+12,5+14+9,5+13}}{\trou{5}} = \displaystyle\frac{\trou{60}}{\trou{5}}=\trou{12}$\\
	La moyenne des notes d'Alice est égale à $12/20$.\\
	
	\vspace{0.05cm}
	Dans cet exemple: 
\includegraphics[width=0.50\linewidth]{don2}
\end{ex}


%%%%%%%%%%%%%%%%%%%%%%%%%%%%%%%%%%%%%%%%%%%%%%%%%%%%%%%%%%%%%%%%%%%%%%%%%%


\section{Représenter graphiquement des données numériques et non numériques}

\begin{prop}{Diagramme en bâtons}{}
	Dans un \textbf{diagramme en bâtons}, la \textbf{hauteur} de chaque bâton est \textbf{proportionnelle} à l'\textbf{effectif} de la valeur qu'il représente.
\end{prop}

\begin{ex}{}{}
	Le professeur de mathématiques a relevé les notes de ses élèves au dernier contrôle:
	\begin{center}
				\includegraphics[width=0.85\linewidth]{don3}
	\end{center}
	Chaque note est représentée par un bâton, la hauteur du bâton correspond à \trou{l'effectif de la note} :
		\begin{center}
		\includegraphics[width=0.5\linewidth]{don4}
	\end{center}
\end{ex}


\begin{prop}{Histogramme}{}
	Quand les données sont nombreuses, on peut les regrouper en \textbf{classes} et les représenter par un \textbf{histogramme}.
\end{prop}

\begin{ex}{}{}
	Lors d'une visite médicale, on a mesuré la taille en centimètres des élèves d'une classe de cinquième. Comme les données sont nombreuses, elles ont été regroupées en classe d'\textbf{amplitude} $5$cm:
		\begin{center}
		\includegraphics[width=0.857\linewidth]{don5} 			\includegraphics[width=0.5\linewidth]{don6}
		\end{center}
\end{ex}




\begin{prop}{Diagramme circulaire}{}
	Dans un \textbf{diagramme circulaire}, les \textbf{mesures des angles sont proportionnelles aux effectifs} de chaque catégorie.
\end{prop}


\begin{ex}{}{}
	Voici la répartition des $100$ élèves de cinquième d'un collège selon leur seconde langue vivante:
			\begin{center}
		\includegraphics[width=0.75\linewidth]{don8}
	\end{center}
L'effectif total est $100$, il correspond à $360 ^{\circ}$ sur le diragramme circulaire.\\
Il suffit donc de multiplier chaque effectif par $3,6$ pour obtenir la mesure de l'angle correspondant.\\
On reporte les résultats obtenues dans le tableau, puis on construit le diagramme.
			\begin{center}
		\includegraphics[width=0.85\linewidth]{don9}
		\includegraphics[width=0.2\linewidth]{don10}
			\end{center}
\end{ex}

\begin{rem}{}{}
	On peut également constuire un diagramme semi-circulaire. La somme des mesures des angles est alors égale à $180 ^{\circ}$.\\
\end{rem}

%%%%%%%%%%%%%%%%%%%%%%%%%%%%%%%%%%%%%%%%%%%%%%%%%%%%%%%%%%%%%%%%%%%%%%%%%%%%%


\section{Carte mentale}

\begin{center}
		\includegraphics[width=0.99\linewidth]{cmdonnees}
\end{center}







%%%%%%%%%%%%%%%%%%%%%%%%%%%%%%%%%%%%%%%%%%%%%%%%%%%%%%%%%%%%%%%%%%%%%%%%%%%%
\end{document}
